% GENERATED FILE. Edit canonical lessons, then run npm run generate:books.
% canonical-source-hash: fnv1a64:ea3532d7eaa645c1
% canonical-lessons: PA-C41-je, PA-C41-je-tan, PA-C41-je-write, PA-C41-jadon, PA-C41-lai

\chapter{If, and When}
\label{ch:pa-if-and-when}
\hlchaptermodality{41}

\begin{chapteropening}
\textbf{By the end of this chapter:} \emph{I can suppose something with a two-part condition, time one event against another, and say what something is for.}

The first joining words that open their clause rather than following it, and the j- family they belong to. Punjabi marks BOTH ends of a condition where English marks only the first. lai is the counterweight: it looks forward to a purpose where kyunki looks back at a cause, and unlike the j- words it FOLLOWS what it applies to, so position alone tells the reader which habit a small word belongs to.
\end{chapteropening}

\section[je]{\pa{ਜੇ} (je) — if}

\label{lesson:PA-C41-je}



\begin{quote}
From the chapter behind you, the word that answers \emph{why}. Two chapters back, the word that sets one clause against another.

\begin{itemize}
\raggedright
  \item \emph{Recall:} \emph{kyoṁki}, then \emph{par}
\end{itemize}

Every joining word so far has taken two things that are \textbf{true}. This one takes something that \textbf{might not be}.
\end{quote}

\begin{cousinweb}
\begin{quote}
\textbf{\pa{ਜੇ}} \emph{(something that might be so)} …
\end{quote}

\textbf{\pa{ਜੇ}} stands at the \textbf{front} of the condition, where English also puts \emph{if}. That is worth noticing: most of the small words in this book — the postpositions, \textbf{\pa{ਮੈਨੂੰ}}'s ending — come \textbf{after} what they apply to. \textbf{\pa{ਜੇ}} comes first, and that alone tells the ear a condition is starting.
\end{cousinweb}

\subsection*{Your turn}
\begin{itemize}
\raggedright
  \item \emph{Say it:} \emph{je}, then any sentence you own.
  \item \emph{Contrast:} the sentence flat, then with \emph{je} in front.
  \item \emph{Build:} suppose something that might not be true.
\end{itemize}

\begin{tcolorbox}[breakable,colback=teal!4,colframe=teal!35!black,title={Before you move on}]
Where does \textbf{\pa{ਜੇ}} stand, and why is that worth noticing? (At the front — most small words in this book come after what they apply to.)
\end{tcolorbox}
\section[je … tāṁ …]{\pa{ਜੇ} … \pa{ਤਾਂ} … (je … tāṁ …) — if … then …}

\label{lesson:PA-C41-je-tan}



\begin{quote}
One lesson back, the word that opens a condition. Four chapters back, where \textbf{\pa{ਨਹੀਂ}} stands.

\begin{itemize}
\raggedright
  \item \emph{Recall:} \emph{je}, and \emph{nahī̃} before the verb
  \item \emph{Recall:} \emph{lekin} — from two chapters back
\end{itemize}
\end{quote}

\begin{cousinweb}
Punjabi usually marks \textbf{both} ends:

\begin{quote}
\textbf{\pa{ਜੇ}} \emph{(condition)}, \textbf{\pa{ਤਾਂ}} \emph{(result)}
\end{quote}

English says \emph{if it is late, I will go} and marks only the front. Punjabi puts \textbf{\pa{ਤਾਂ}} at the head of the result as well, so the sentence announces its two halves. Beginners often leave it out; it is understood, and the full pair is what you will hear.

The result half is free to be negative, which is where this meets the chapter that opened this run.
\end{cousinweb}

\subsection*{Your turn}
\begin{itemize}
\raggedright
  \item \emph{Build:} \emph{je} + a condition, \emph{tāṁ} + a result.
  \item \emph{Build:} the same, with the result negative.
  \item \emph{Contrast:} the pair with \emph{tāṁ}, then without it.
\end{itemize}

\begin{tcolorbox}[breakable,colback=teal!4,colframe=teal!35!black,title={Before you move on}]
Make a two-part condition. Which half does English leave unmarked that Punjabi marks? (The result — \textbf{\pa{ਤਾਂ}}.)
\end{tcolorbox}
\section[je]{\pa{ਜੇ} reaches the page}

\label{lesson:PA-C41-je-write}



\begin{quote}
One chapter back you wrote \textbf{\pa{ਕਿ}} and \textbf{\pa{ਕੀ}} and checked which side each stroke fell on. Write \textbf{\pa{ਕਿ}} once from memory.
\end{quote}

\subsection*{Script — two signs, neither new}
\begin{quote}
\textbf{\pa{ਜ}} — \emph{ja} · \textbf{\pa{ੇ}} — the laav stroke above
\end{quote}

Set \textbf{\pa{ਅਤੇ}}, \textbf{\pa{ਜੇ}} side by side and look only at the top. It is the same laav in both, sitting on different letters. One stroke has now spelled several of this book's most useful words.

\subsection*{Writing — guided copy, model in view}
Copy \textbf{\pa{ਜੇ}} once with the model in view, then write it at the front of a line, which is where it belongs.

\begin{tcolorbox}[breakable,colback=teal!4,colframe=teal!35!black,title={Before you move on}]
Write \textbf{\pa{ਜੇ}} from memory. Which stroke does it share with \textbf{\pa{ਅਤੇ}}? (The laav above.)

\begin{itemize}
\raggedright
  \item \emph{Recall:} \textbf{\pa{ਪਰ}} written — from two chapters back
\end{itemize}
\end{tcolorbox}
\section[jadoṁ]{\pa{ਜਦੋਂ} (jadoṁ) — when}

\label{lesson:PA-C41-jadon}



\begin{quote}
One lesson back you wrote \textbf{\pa{ਜੇ}}. And before that, the two-part condition.

\begin{itemize}
\raggedright
  \item \emph{Recall:} \textbf{\pa{ਜੇ}}, and \emph{je … tāṁ …}
  \item \emph{Recall:} the pair rule — from two chapters back
\end{itemize}
\end{quote}

\begin{cousinweb}
\begin{quote}
\textbf{\pa{ਜਦੋਂ}} \emph{(one event)}, \emph{(the other)}
\end{quote}

\textbf{\pa{ਜਦੋਂ}} sets two events against each other in time. Like \textbf{\pa{ਜੇ}}, it opens its clause rather than following it — and look at the front of both words. The same \textbf{\pa{ਜ}}.

That is not a coincidence. Punjabi builds a family of clause-openers on that one letter, the way it builds its questions on \textbf{\pa{ਕ}}. You now own two members of it, and the pattern is worth more than either word alone.
\end{cousinweb}

\subsection*{Your turn}
\begin{itemize}
\raggedright
  \item \emph{Build:} \emph{jadoṁ} + an event, then a second event.
  \item \emph{Contrast:} \emph{je} and \emph{jadoṁ} on the same two clauses — one supposes, one times.
  \item \emph{Say it:} the shared \textbf{\pa{ਜ}} at the front of both.
\end{itemize}

\begin{tcolorbox}[breakable,colback=teal!4,colframe=teal!35!black,title={Before you move on}]
Say a \emph{when} sentence. What do \textbf{\pa{ਜੇ}} and \textbf{\pa{ਜਦੋਂ}} share, and what does it mark? (The opening \textbf{\pa{ਜ}} — a family of clause-openers.)
\end{tcolorbox}
\section[laī]{\pa{ਲਈ} (laī) — for, in order to}

\label{lesson:PA-C41-lai}



\begin{quote}
One lesson back, the word that times one event against another. One chapter back, the word that answers \emph{why}.

\begin{itemize}
\raggedright
  \item \emph{Recall:} \emph{jadoṁ}, and \emph{kyoṁki}
  \item \emph{Recall:} the contrast kit — from two chapters back
\end{itemize}
\end{quote}

\begin{cousinweb}
\textbf{\pa{ਕਿਉਂਕਿ}} looks \textbf{backwards} — it names the cause of something that already happened. \textbf{\pa{ਲਈ}} looks \textbf{forwards} — it names what something is \textbf{for}:

\begin{quote}
\emph{(a thing)} \textbf{\pa{ਲਈ}} — \textquotedblleft{}for …\textquotedblright{}
\end{quote}

And notice where it stands. \textbf{\pa{ਜੇ}} and \textbf{\pa{ਜਦੋਂ}} open their clauses; \textbf{\pa{ਲਈ}} \textbf{follows} what it applies to, like the postpositions this book taught long ago. Two habits, and the word tells you which one it belongs to by where it sits.

Cause and purpose are different answers to the same \emph{why}, and a language that used one word for both would make its learners guess. Punjabi does not.
\end{cousinweb}

\subsection*{Your turn}
\begin{itemize}
\raggedright
  \item \emph{Say it:} something you can name, then \emph{laī}.
  \item \emph{Contrast:} a \emph{kyoṁki} answer and a \emph{laī} answer to the same question.
  \item \emph{Contrast:} \emph{je} at the front, \emph{laī} at the back.
\end{itemize}

\begin{tcolorbox}[breakable,colback=teal!4,colframe=teal!35!black,title={Before you move on}]
Which of \textbf{\pa{ਕਿਉਂਕਿ}} and \textbf{\pa{ਲਈ}} looks back at a cause, and which forward to a purpose? Where does each stand? Next chapter: \textbf{he}, \textbf{we}, and the word for \emph{the one who}.
\end{tcolorbox}
